\section*{Chapitre \ref{ch:g11:dissolution} --- La dissolution des solides ioniques et moléculaires}

\begin{solution}{exo:g11:dissolution:1}
\ce{NaCl(s) -> Na+(aq) + Cl-(aq)};
\ce{CaCl2(s) -> Ca^{2+}(aq) + 2Cl-(aq)};
\ce{Na2SO4(s) -> 2Na+(aq) + SO4^{2-}(aq)};
\ce{AlCl3(s) -> Al^{3+}(aq) + 3Cl-(aq)}.
\end{solution}

\begin{solution}{exo:g11:dissolution:2}
$[\ce{Na+}] = 2 \times 0.050 = \qty{0.10}{mol/L}$;
$[\ce{SO4^{2-}}] = \qty{0.050}{mol/L}$.
\end{solution}

\begin{solution}{exo:g11:dissolution:3}
L'\omterm{def:g11:dissolution:solvation}{hydratation} est le fait que l'\omterm{def:g9:ions:ion}{ion} s'entoure de \omterm{def:g7:atoms-and-molecules:molecule}{molécules} d'eau
retenues par lui. L'extrémité oxygène ($\delta^-$) est tournée vers un
\omterm{def:g9:ions:ion}{cation}, l'extrémité hydrogène ($\delta^+$) vers un \omterm{def:g9:ions:ion}{anion}~: les charges
opposées s'attirent.
\end{solution}

\begin{solution}{exo:g11:dissolution:4}
$M(\ce{AlCl3}) = 27.0 + 3 \times 35.5 = \qty{133.5}{g/mol}$;
$n = 2.67 / 133.5 = \qty{0.0200}{mol}$; $c = 0.0200 / 0.2000 =
\qty{0.100}{mol/L}$; $[\ce{Al^{3+}}] = \qty{0.100}{mol/L}$,
$[\ce{Cl-}] = \qty{0.300}{mol/L}$.
\end{solution}

\begin{solution}{exo:g11:dissolution:5}
La tête \ce{-COO-} est ionique, attirée par l'eau~: elle est \omterm{def:g11:dissolution:hydrophilic}{hydrophile}.
La queue de 17 \omterm{def:g7:atoms-and-molecules:atom}{atomes} de carbone, qui ne comporte que des liaisons
\ce{C-C} et des liaisons \ce{C-H} presque non polarisées, est apolaire~:
elle est \omterm{def:g11:dissolution:hydrophilic}{hydrophobe}.
\end{solution}

\begin{solution}{exo:g11:dissolution:6}
Le sel~: l'eau (solide ionique, \omterm{def:g6:solutions-and-solubility:solute}{solvant} polaire). La cire~: le
cyclohexane (chaînes apolaires). Le sucre~: l'eau (ses groupes \ce{O-H}
forment des \omterm{def:g11:polarity-and-cohesion:hydrogen-bond}{liaisons hydrogène} avec l'eau). Le diiode~: le cyclohexane
(\omterm{def:g11:polarity-and-cohesion:polar-molecule}{molécule apolaire}).
\end{solution}

\begin{solution}{exo:g11:dissolution:7}
Le groupe \ce{O-H} de l'éthanol forme avec les \omterm{def:g7:atoms-and-molecules:molecule}{molécules} d'eau des
\omterm{def:g11:polarity-and-cohesion:hydrogen-bond}{liaisons hydrogène} qui remplacent celles qu'il rompt~: l'éthanol trouve
sa place dans l'eau. L'hexane, apolaire, ne peut former aucune liaison
de ce type~; les \omterm{def:g7:atoms-and-molecules:molecule}{molécules} d'eau, retenues les unes aux autres par des
\omterm{def:g11:polarity-and-cohesion:hydrogen-bond}{liaisons hydrogène}, ne le laissent pas entrer.
\end{solution}

\begin{solution}{exo:g11:dissolution:8}
Les queues \omterm{def:g11:dissolution:hydrophilic}{hydrophobes} fuient l'eau et se dissolvent dans la graisse, au
centre~; les têtes chargées sont attirées par l'eau et restent à
l'extérieur. Toutes les \omterm{def:g11:dissolution:amphiphilic}{micelles} portent des charges négatives à leur
surface, et des charges de même signe se repoussent~: elles restent
séparées.
\end{solution}

\begin{solution}{exo:g11:dissolution:9}
Le cyclohexane~: il n'est pas \omterm{def:g6:solutions-and-solubility:miscible}{miscible} avec l'eau, si bien qu'il forme
une phase séparée que l'on peut évacuer, et le parfum (apolaire) s'y
dissout bien. L'éthanol est \omterm{def:g6:solutions-and-solubility:miscible}{miscible} avec l'eau~: aucune seconde phase
ne se forme, et rien ne peut être séparé.
\end{solution}

\begin{solution}{exo:g11:dissolution:10}
Le cyclohexane est moins dense que l'eau. C'est la phase inférieure,
aqueuse, que l'on évacue par le robinet. Le diiode se trouve dans la
phase supérieure de cyclohexane, qui reste dans l'ampoule.
\end{solution}

\begin{solution}{exo:g11:dissolution:11}
Rien~: le sulfate de cuivre est ionique et ne \omterm{def:g2:mixing-and-dissolving:dissolve}{se dissout} pas dans le
cyclohexane, apolaire. L'eau reste bleue et le cyclohexane incolore.
\end{solution}

\begin{solution}{exo:g11:dissolution:12}
$[\ce{Cl-}] = 2c$, so $c = \qty{0.150}{mol/L}$;
$n = 0.150 \times 0.2500 = \qty{0.0375}{mol}$;
$m = 0.0375 \times 111.1 = \qty{4.17}{g}$.
\end{solution}

\begin{solution}{exo:g11:dissolution:13}
Les \omterm{def:g9:ions:ion}{ions} calcium et magnésium précipitent les \omterm{def:g9:ions:ion}{ions} stéarate sous forme
de dépôt~: le savon est consommé avant de pouvoir former des \omterm{def:g11:dissolution:amphiphilic}{micelles}, et
il mousse mal. Le calcium seul~: \qty{0.010}{mol} de \ce{Ca^{2+}}
accapare \qty{0.020}{mol} de stéarate, soit $0.020 \times 306.0 =
\qty{6.1}{g}$ de stéarate de sodium par litre.
\end{solution}

\begin{solution}{exo:g11:dissolution:14}
Volume \qty{0.2000}{L}. \ce{Na+}~: $0.020 / 0.2000 = \qty{0.10}{mol/L}$~;
\ce{Ca^{2+}}~: $0.010 / 0.2000 = \qty{0.050}{mol/L}$~; \ce{Cl-}~:
$(0.020 + 0.020) / 0.2000 = \qty{0.20}{mol/L}$. Charge positive~:
$0.10 + 2 \times 0.050 = 0.20$~; négative~: $0.20$. Le \omterm{def:g6:pure-substances-and-mixtures:mixture}{mélange} est
neutre.
\end{solution}

\begin{solution}{exo:g11:dissolution:15}
Un litre d'heptane a une masse de \qty{0.68}{kg} et peut contenir
$17 \times 0.68 = \qty{12}{g}$ de diiode, environ 40 fois plus qu'un
litre d'eau. Agités ensemble, presque tout le diiode passe dans
l'heptane~: l'\omterm{def:g10:chemical-species:extraction}{extraction} est efficace.
\end{solution}

\begin{solution}{pb:g11:dissolution:1}
\textbf{1.} Des roches que l'eau a traversées (calcaire, dolomie,
gypse), qui se dissolvent un peu.

\textbf{2.} $[\ce{Ca^{2+}}] = 0.120 / 40.1 = \qty{2.99e-3}{mol/L}$~;
$[\ce{Mg^{2+}}] = 0.024 / 24.3 = \qty{9.9e-4}{mol/L}$.

\textbf{3.} \ce{Ca(HCO3)2 -> Ca^{2+}(aq) + 2HCO3-(aq)}~; donc
$[\ce{HCO3-}] = 2 \times \num{2.99e-3} = \qty{5.99e-3}{mol/L}$.

\textbf{4.} $\num{2.99e-3} \times 150 = \qty{0.449}{mol}$.

\textbf{5.} $18 \times 12.0 + 35 \times 1.0 + 2 \times 16.0 + 23.0 =
\qty{306.0}{g/mol}$.

\textbf{6.} \ce{C17H35COONa(s) -> C17H35COO-(aq) + Na+(aq)}.

\textbf{7.} La tête \ce{-COO-}~: \omterm{def:g11:dissolution:hydrophilic}{hydrophile}~; la queue \ce{C17H35-}~:
\omterm{def:g11:dissolution:hydrophilic}{hydrophobe}. Il possède les deux sortes de parties~: il est \omterm{def:g11:dissolution:amphiphilic}{amphiphile}.

\textbf{8.} Une minuscule sphère d'\omterm{def:g9:ions:ion}{ions} de savon, queues vers
l'intérieur, têtes vers l'extérieur. Les queues se dissolvent dans la
graisse, qui se divise en gouttelettes enveloppées de \omterm{def:g11:dissolution:amphiphilic}{micelles}~; les
surfaces chargées maintiennent les gouttelettes séparées dans l'eau, qui
les emporte au rinçage.

\textbf{9.} Le savon doit former des \omterm{def:g9:ions:ion}{ions} et des \omterm{def:g11:dissolution:amphiphilic}{micelles}, ce qui
demande de l'eau autour des têtes~; et la graisse est emportée par l'eau
de rinçage. Dans l'huile, la graisse resterait simplement dissoute sur
la peau.

\textbf{10.} \ce{2C17H35COO- + Ca^{2+} -> Ca(C17H35COO)2}~: charges
$2 \times (-1) + 2 = 0$ à gauche, $0$ à droite.

\textbf{11.} \ce{C17H35COO-}~: excès $- 2x$~; \ce{Ca^{2+}}~:
$\num{2.99e-3} - x$~; stéarate de calcium~: $x$. Les \omterm{def:g9:ions:ion}{ions} calcium sont
limitants~: $x_{\max} = \qty{2.99e-3}{mol}$.

\textbf{12.} $2 \times \num{2.99e-3} = \qty{5.99e-3}{mol}$ par litre.

\textbf{13.} $M = 36 \times 12.0 + 70 \times 1.0 + 4 \times 16.0 + 40.1 =
\qty{606.1}{g/mol}$~; $\num{2.99e-3} \times 606.1 = \qty{1.81}{g}$ de
dépôt par litre.

\textbf{14.} $2 \times \num{9.9e-4} = \qty{1.98e-3}{mol}$ par litre.

\textbf{15.} $2 \times 0.449 = \qty{0.898}{mol}$.

\textbf{16.} $0.898 \times 306.0 = \qty{275}{g}$.

\textbf{17.} Le magnésium~: $\num{9.9e-4} \times 150 = \qty{0.148}{mol}$,
qui accapare \qty{0.296}{mol} de stéarate, soit \qty{91}{g} de savon~; au
total, environ $275 + 91 = \qty{366}{g}$.

\textbf{18.} $275 / 100 \approx 2.7$ pains.

\textbf{19.} Le calcium d'un bain gaspille environ \qty{0.27}{kg} de
savon en dépôt.
\end{solution}
